\chapter{Модель эмбеддингов GigaEmbeddings}\label{ch:gigaembeddings}

Текстовые эмбеддинги --- векторные представления, кодирующие семантическую информацию естественного языка, --- являются фундаментальным компонентом широкого спектра приложений NLP: информационного поиска, вопросно-ответных систем, оценки семантического сходства, битекст-майнинга (добычи параллельных текстов) и рекомендательных систем. В конвейерах информационного поиска эмбеддинги обеспечивают эффективный первичный этап извлечения через приближённый поиск ближайших соседей, а в системах генерации с извлечением информации (RAG)~\cite{2005.11401} --- динамическое заземление больших языковых моделей на внешних знаниях.

В настоящей главе представлен фреймворк GigaEmbeddings для обучения высокопроизводительных русскоязычных текстовых эмбеддингов на основе иерархического инструктивного дообучения декодерной LLM GigaChat-3B, являющейся частью семейства GigaChat~\cite{2506.09440}, описанного в Главе~\ref{ch:gigachat}. Результаты настоящей главы опубликованы в~\cite{2510.22369}.

%% =================================================================
\section{Архитектура}\label{sec:ge_architecture}
%% =================================================================

\subsection{Базовая модель}\label{subsec:ge_backbone}

В качестве основы GigaEmbeddings используется предобученный чекпоинт GigaChat-3B --- декодерная LLM семейства GigaChat. Выбор модели масштаба 3B параметров отражает баланс между вычислительной эффективностью и репрезентативной мощностью модели: при ограниченном доступе к крупномасштабным GPU-кластерам (8$\times$A100 80GB) приоритет отдан масштабу, обеспечивающему эффективное дообучение и инференс при сохранении достаточной выразительности для сложных задач русского языка.

\subsection{Двунаправленное внимание и прунинг слоёв}\label{subsec:ge_bidirectional}

Каузальные маски внимания удалены на этапе обучения для обеспечения двунаправленного контекстного моделирования, что позволяет каждому токену учитывать информацию как от предшествующих, так и от последующих позиций. Данная модификация критически важна для задач построения эмбеддингов, где полный контекст последовательности необходим для формирования качественного векторного представления.

Следуя результатам Gromov et al.~\cite{gromov2025unreasonable}, выполнен прунинг 25\% глубоких слоёв трансформера: 9 из 36 блоков (модули самовнимания и feed-forward) удалены. В результате в качестве основы модели эмбеддингов используется LLM на 2.5B параметров, что обеспечивает значительное снижение задержки инференса при пренебрежимом снижении качества (см. Раздел~\ref{subsec:ge_pruning_ablation}).

\subsection{Латентный пулинг внимания}\label{subsec:ge_pooling}

Для агрегации скрытых активаций декодера в финальный вектор эмбеддинга используется голова латентного пулинга внимания (latent attention pooling), вдохновлённая подходом NV-Embed~\cite{2405.17428}. Архитектура модели представлена на Рисунке~\ref{fig:ge_arch}.

\begin{figure}[H]
    \centering
    \includegraphics[width=0.85\textwidth]{data/gigaembeddings-efficient-russian-language-embedding-model/figure-3-1.jpg}
    \caption{Архитектура GigaEmbeddings. Модель состоит из LLM GigaChat и слоя латентного внимания. Слой работает через механизм кросс-внимания: выход декодера выступает в роли запроса~($\mathbf{Q}$), а обучаемый латентный массив предоставляет ключ~($\mathbf{K}$) и значение~($\mathbf{V}$). Выход внимания обрабатывается многослойным перцептроном (MLP).}
    \label{fig:ge_arch}
\end{figure}

%% =================================================================
\section{Конвейер обучения}\label{sec:ge_training}
%% =================================================================

Предложен иерархический трёхэтапный конвейер инструктивного дообучения, последовательно уточняющий эмбеддинги для разнообразных прикладных задач.

\subsection{Функция потерь}\label{subsec:ge_loss}

Используется классическая функция потерь InfoNCE~\cite{oord2019representation} с фиксированной температурой $\tau = 0.02$:
\begin{equation}
\label{eq:infonce}
    \mathcal{L} = -\log \frac{\varphi(q, d^+)}{\varphi(q, d^+) + \sum_{n_i \in \mathcal{N}} \varphi(q, n_i)},
\end{equation}
где $\mathcal{N}$ обозначает множество всех внутрибатчевых, кросс-батчевых и жёстких негативных примеров, а $\varphi(q, d^+)$ --- функция, вычисляющая оценку совпадения между запросом $q$ и позитивным документом $d^+$:
\begin{equation}
\label{eq:matching}
    \varphi(q, d^+) = \exp\!\left(\frac{1}{\tau} \cos(\mathbf{h}_q, \mathbf{h}_{d^+})\right),
\end{equation}
где $\mathbf{h}_q, \mathbf{h}_{d^+}$ --- векторные эмбеддинги, полученные из модели.

\subsection{Этап 1: Контрастивное предобучение}\label{subsec:ge_stage1}

Первый этап использует корпуса веб-масштаба в формате <<заголовок--пассаж>>, извлечённые из Wikipedia, Reddit, StackExchange и S2ORC~\cite{lo2020s2orc}, а также добытые текстовые пассажи. Для повышения разнообразия запросов применяется инструктивно-дообученная LLM для синтетической генерации контекстуально релевантных запросов к каждому пассажу.

Обучение использует потери InfoNCE с внутрибатчевыми негативами (из того же батча) и кросс-батчевыми негативами (кэшированными из недавних батчей), обеспечивая эффективное контрастивное обучение на 16\,384 образцах в батче. Большой размер батча максимизирует разнообразие негативных примеров при адаптации декодерной архитектуры GigaChat к двунаправленному вниманию для задач построения эмбеддингов.

\subsection{Этап 2: Дообучение с жёсткими негативами}\label{subsec:ge_stage2}

Второй этап фокусируется на высококачественных размеченных датасетах, включая MS-MARCO~\cite{bajaj2018msmarco}, Natural Questions (NQ)~\cite{kwiatkowski2019nq}, SQuAD~\cite{rajpurkar2016squad}, MIRACL~\cite{zhang2023miracl} и Mr~TyDi~\cite{zhang2021mrtydi}, для уточнения способностей к информационному поиску. Функция потерь InfoNCE расширена включением 7 курируемых жёстких негативных примеров на запрос, добытых через пороги семантического сходства, наряду с внутрибатчевыми негативами. Размер батча снижен до 512 образцов для приоритизации точности над масштабом, обеспечивая робастное дискриминативное обучение для сложных сценариев поиска.

\subsection{Этап 3: Мультизадачная генерализация}\label{subsec:ge_stage3}

Финальный этап вводит задачи классификации и кластеризации в обучающую смесь задач информационного поиска, классификации и кластеризации для расширения применимости модели. Для предотвращения ложных негативов в нерелевантных задачах удалены внутрибатчевые негативы с задачно-специфичным инструктивным дообучением, с использованием унифицированной функции потерь InfoNCE для всех целей. Размер батча 512 образцов обеспечивает баланс между вычислительной эффективностью и качеством на бенчмарках классификации и кластеризации.

Методология использует стратегию выборки на основе примеров: позитивные экземпляры выбираются из того же класса/кластера, что и запрос, а негативные --- из других классов/кластеров, обеспечивая внутриклассовую когезию и межклассовое различие при обучении.

\subsection{Гиперпараметры обучения}\label{subsec:ge_hyperparams}

Гиперпараметры для всех трёх этапов обучения представлены в Таблице~\ref{tab:ge_hyperparams}.

\begin{figure}[H]
    \centering
    \includegraphics[width=0.7\textwidth]{data/gigaembeddings-efficient-russian-language-embedding-model/table-8-5.jpg}
    \caption{Гиперпараметры для этапов контрастивного предобучения, дообучки и мультизадачного обучения.}
    \label{tab:ge_hyperparams}
\end{figure}

%% =================================================================
\section{Стратегия промптирования}\label{sec:ge_prompting}
%% =================================================================

Поскольку задачи построения эмбеддингов включают как симметричные (классификация, кластеризация, STS), так и асимметричные (поиск, переранжирование) типы, промптирование модели может существенно влиять на качество. Сравнены два способа промптирования: \textit{prefix}~--- добавление фиксированного префикса к входным данным, и \textit{instruct}~--- использование полной инструкции, описывающей задачу.

Результаты сравнения представлены в Таблице~\ref{tab:ge_prompting}. Инструктивное промптирование демонстрирует превосходство на 0.8 балла ruMTEB по сравнению с префиксным подходом.

\begin{figure}[H]
    \centering
    \includegraphics[width=0.5\textwidth]{data/gigaembeddings-efficient-russian-language-embedding-model/table-5-3.jpg}
    \caption{Сравнение способов промптирования: prefix и instruct.}
    \label{tab:ge_prompting}
\end{figure}

%% =================================================================
\section{Экспериментальная оценка}\label{sec:ge_evaluation}
%% =================================================================

\subsection{Базовые линии}\label{subsec:ge_baselines}

Модель сравнивается с рядом сильных базовых линий на бенчмарке MTEB(rus, v1.1)~\cite{snegirev2025rumteb}:
\begin{itemize}
    \item Qwen3-Embedding-8B~\cite{2506.05176};
    \item FRIDA;
    \item multilingual-E5-large-instruct~\cite{2212.03533};
    \item E5-Mistral-7B-instruct~\cite{2401.00368};
    \item GritLM-7B~\cite{muennighoff2025gritlm};
    \item BGE-M3~\cite{2309.07597}.
\end{itemize}

\subsection{Основные результаты}\label{subsec:ge_main_results}

Эффективность модели оценена на русскоязычном бенчмарке MTEB(rus, v1.1) --- актуальной версии ruMTEB~\cite{snegirev2025rumteb}, охватывающей семь категорий задач. По данным публичного рейтинга по состоянию на июнь 2026~г. модель Giga-Embeddings-instruct занимает первое место со средним баллом $74.16$ (Mean over tasks), опережая сильные многоязычные и русскоязычные базовые линии, включая Qwen3-Embedding-8B ($71.42$) и FRIDA ($70.95$). Результаты по категориям приведены в Таблице~\ref{tab:ge_results}.

Модель показывает лучшие результаты в шести из семи категорий: классификации ($76.66$), мультиклассовой классификации ($57.07$), классификации пар ($79.57$), переранжировании ($74.02$) и --- что особенно важно для двухэтапного поиска --- информационном поиске (retrieval, $81.41$), где она опережает Qwen3-Embedding-8B ($78.19$) при значительно меньшем числе параметров. На кластеризации модель практически не уступает лучшей базовой линии ($67.48$ против $67.49$ у FRIDA), а на семантическом текстовом сходстве (STS, $75.75$) уступает Qwen3-Embedding-8B ($79.28$).

\begin{table}[H]
    \centering
    \caption{Сравнение GigaEmbeddings с базовыми линиями на бенчмарке MTEB(rus, v1.1) по семи категориям задач (по данным публичного рейтинга, июнь 2026~г.). Лучший результат в каждом столбце выделен жирным. Обозначения: Cls~--- классификация, Clust~--- кластеризация, MLCls~--- мультиклассовая классификация, PairCls~--- классификация пар, Rerank~--- переранжирование, Retr~--- информационный поиск.}
    \label{tab:ge_results}
    \resizebox{\textwidth}{!}{
    \begin{tabular}{lcccccccc}
        \toprule
        Модель & Среднее & Cls & Clust & MLCls & PairCls & Rerank & Retr & STS \\
        \midrule
        Giga-Embeddings-instruct       & \textbf{74.16} & \textbf{76.66} & 67.48 & \textbf{57.07} & \textbf{79.57} & \textbf{74.02} & \textbf{81.41} & 75.75 \\
        Qwen3-Embedding-8B             & 71.42 & 74.35 & 66.72 & 44.58 & 71.33 & 70.19 & 78.19 & \textbf{79.28} \\
        FRIDA                          & 70.95 & 73.71 & \textbf{67.49} & 52.23 & 66.45 & 71.50 & 76.38 & 74.35 \\
        E5-Mistral-7B-instruct         & 67.04 & 69.07 & 64.24 & 42.93 & 60.81 & 69.96 & 73.29 & 73.71 \\
        GritLM-7B                      & 65.02 & 69.92 & 64.30 & 41.96 & 58.93 & 69.99 & 55.52 & 74.63 \\
        multilingual-E5-large-instruct & 62.41 & 66.28 & 63.13 & 41.15 & 63.89 & 64.35 & 48.45 & 76.48 \\
        BGE-M3                         & 61.63 & 60.44 & 52.38 & 34.86 & 60.60 & 69.71 & 75.20 & 73.68 \\
        \bottomrule
    \end{tabular}
    }
\end{table}

\subsection{Абляционные исследования: прунинг}\label{subsec:ge_pruning_ablation}

Результаты сравнения полной и обрезанной версий модели представлены в Таблице~\ref{tab:ge_pruning}. Прунинг 25\% слоёв приводит к снижению среднего балла на исходной версии бенчмарка (ruMTEB) лишь на 0.2 (с 69.3 до 69.1), что подтверждает эффективность подхода.

\begin{figure}[H]
    \centering
    \includegraphics[width=0.45\textwidth]{data/gigaembeddings-efficient-russian-language-embedding-model/table-5-2.jpg}
    \caption{Сравнение полной и обрезанной версий модели GigaEmbeddings.}
    \label{tab:ge_pruning}
\end{figure}

\subsection{Абляционные исследования: влияние предобучения}\label{subsec:ge_pretrain_ablation}

Для оценки необходимости этапа контрастивного предобучения сравнены два варианта: трёхэтапное обучение (с предобучением) и двухэтапное (только дообучка и мультизадачное обучение). Результаты представлены в Таблице~\ref{tab:ge_pretrain_ablation}.

Включение контрастивного предобучения на слабо размеченных данных дополнительно повышает средний балл на 0.6 (на исходной версии бенчмарка ruMTEB), прежде всего за счёт задач информационного поиска. Данный эффект объясняется необходимостью модели переконфигурироваться для новой задачи и адаптироваться к двунаправленному механизму внимания.

\begin{figure}[H]
    \centering
    \includegraphics[width=0.55\textwidth]{data/gigaembeddings-efficient-russian-language-embedding-model/table-6-4.jpg}
    \caption{Влияние этапа контрастивного предобучения на качество модели.}
    \label{tab:ge_pretrain_ablation}
\end{figure}

%% =================================================================
\section{Выводы}\label{sec:ge_conclusions}
%% =================================================================

В настоящей главе представлен GigaEmbeddings --- трёхэтапный фреймворк инструктивного дообучения для обучения высокопроизводительных текстовых эмбеддингов на основе декодерной LLM GigaChat-3B. Иерархический конвейер, включающий масштабное контрастивное предобучение, дообучку с жёсткими негативами и мультизадачную генерализацию, решает ключевые ограничения существующих методов, объединяя задачи информационного поиска, классификации и кластеризации с использованием синтетической генерации данных.

Оценка на бенчмарке MTEB(rus, v1.1) демонстрирует результаты уровня state-of-the-art: модель занимает первое место со средним баллом 74.16, превосходя сильные базовые линии с большим числом параметров, такие как Qwen3-Embedding-8B и E5-Mistral-7B. Абляционные исследования подтверждают необходимость контрастивного предобучения (прирост 0.6 балла) и превосходство инструктивного промптирования над префиксным подходом. Стратегический прунинг на основе результатов Gromov et al.~\cite{gromov2025unreasonable} сохраняет полную производительность модели при значительном повышении эффективности.

Модель выпущена в открытый доступ под лицензией MIT (Giga-Embeddings-instruct~\cite{gigaembeddings_hf}), обеспечивая воспроизводимую основу для дальнейших исследований. Работа демонстрирует жизнеспособность декодерных LLM в качестве универсальных энкодеров и потенциал парадигм обучения на синтетических данных.
